"\\pagestyle{plain}\n\\chapter*{Agradecimientos}\n\nA mis tutores, Manuel y Tomás, por su inestimable guía, su constante disponibilidad y el rigor metodológico que han aportado a la dirección de este trabajo. Sus valiosas orientaciones críticas y su confianza en mí a lo largo de este proyecto han sido fundamentales para superar los desafíos técnicos y conceptuales que plantea la optimización cuántica.\n\n\\vspace{1cm}\nAl Centro de Supercomputación de Galicia (CESGA), con especial mención a Carlos Fernández Sánchez, por su asistencia y por facilitar el acceso a las capacidades de supercomputación del FinisTerrae III y a la QPU superconductora de Qmio. La posibilidad de validar experimentalmente este modelo en hardware cuántico real ha sido un pilar indispensable y una oportunidad excepcional en mi formación.\n\n\\vspace{1cm}\nA mi familia, pareja y amigos, por su comprensión incondicional, su apoyo constante y por ser el refugio necesario durante las exigentes jornadas de este trabajo. Gracias por hacer este camino académico mucho más ameno, compartido y enriquecedor.\n"